\documentclass[10pt,a4paper]{amsart}
\usepackage[margin=19mm]{geometry}
\usepackage{amsmath,amssymb,mathtools}
\usepackage[scr=rsfs,cal=euler]{mathalfa}
\usepackage{xfrac}
\usepackage[unicode,hidelinks,bookmarks=false]{hyperref}
\newcommand{\ie}{i.e.\ }
\newcommand{\cf}{cf.\ }
\newcommand{\resp}{respectively\ }
\newcommand{\Subsection}{\subsection}
\providecommand{\nobreakdash}{\nobreak\hbox{-}}
\setlength{\emergencystretch}{2em}
\multlinegap0pt
\usepackage{amssymb}
\usepackage{mathtools}
\usepackage[shortlabels]{enumitem}
\usepackage{tikz}
\newtheorem{theorem}{Theorem}
\title{A Manual for Ends, Semistability and Simple Connectivity at Infinity for Groups and Spaces}
\author{Michael Mihalik}
\date{Source: arXiv:2507.17060v2 (CC BY 4.0)}
\begin{document}
\maketitle

\noindent\emph{Source and licence.} A Manual for Ends, Semistability and Simple Connectivity at Infinity for Groups and Spaces. Version arXiv:2507.17060v2 (CC BY 4.0), distributed by arXiv under the Creative Commons Attribution 4.0 International licence, http://creativecommons.org/licenses/by/4.0/, as recorded in the arXiv licence metadata for this exact version and shown on the abstract page https://arxiv.org/abs/2507.17060v2. Full licence notice: \texttt{licenses/arxiv-2507.17060-CC-BY-4.0.txt}.\par\smallskip

\noindent\textit{Editorial scope.} Counted unit: the article's theorem theorem (source0.tex lines 582--596). The article's complete original proof of it is delivered with it (source0.tex lines 597--702, one \texttt{proof} environment of 106 lines). No mathematics is written, repaired or re-derived: the statement and the proof are the article's own lines, in the article's own notation, and no retained line is substituted or re-flowed.  Retained uncounted as the source-local context the proof needs: lines 557--559.  One declared adaptation: exactly 2 ancillary strings inside the retained spans are omitted (image-inclusion commands and, where the source repeats a label, the repeated label definitions) and no other line differs from the source; the figure environments, with their placement, captions and labels, are kept, and the package records the omitted strings in fidelity receipts bound to the affected bodies.  No retained line contains a citation, so the wrapper carries no bibliography and no bibliography entry.  The retained lines contain the article's own diagram code, which the wrapper typesets with the standard tikz packages; no image file is delivered and the package declares no asset dependency.  Editorial formatting only, in addition: an AMS \texttt{amsart} preamble that restates the article's own declarations and macros used by the retained lines, namely the environments \texttt{theorem} on separate counters, together with the standard packages those lines need.  The retained environments print in this document as Theorem 1 in order of appearance, whereas the article prints the same items with the numbering of its own full document.  The complete native source of the article as distributed in the arXiv e-print for this exact version is bundled unchanged as \texttt{sources/182053/source0.tex}.
\begin{theorem} \label{OneEndSplit} 
Suppose $G$ is a finitely generated group and $G$ splits as $A\ast_C B$ or $A\ast_C$ where $C$ is not equal to $A$ or $B$ and $C$ is finite. If $H$ is 1-ended in $G$, then $H$ is a subgroup of a conjugate of $A$ or $B$. 
\end{theorem}

\begin{theorem}
Suppose $G$ is a finitely generated group that splits non-trivially as  $A\ast_C B$ or $A\ast_C$ where $C$ is not equal to $A$ or $B$ and $C$ is finite. Let $\mathcal T$ be the Bass-Serre tree for this splitting and suppose that the subgroup  $H$ is 1-ended in  $G$. Let $\mathcal T_1$ be the subtree of $\mathcal T$ containing all vertices $hA$ and $hB$ for $h\in H$, so that $H$ acts on $\mathcal T_1$ by graph isomorphisms. Theorem \ref{OneEndSplit} implies $H$ is a subgroup of an $A$ or $B$ coset $Q$ of $G$. Then

(1) If $D$ is the distance in $\mathcal T$ from $A$ to $Q$, then for all $h\in H$, the distance from $hA$ to $Q$ is $D$. Similarly for $B$. 

(2) If $h_1, h_2\in H$ are such that $h_1A\ne h_2A$ and $\alpha$ is the (even length) geodesic in $\mathcal T_1$ between $h_1A$ and $h_2A$, then the closest point $M$ of $\alpha$ to $Q$  is the midpoint (vertex) of $\alpha$. The geodesics from $h_1A$ to $Q$ and $h_2A$ to $Q$, both contain $M$. Furthermore, $M$ is stabilized by $h_2h_1^{-1}$. In particular, if $E$ is the first edge of $\alpha$, then the last edge of $\alpha$ is $h_2h_1^{-1}(E)$. Similarly for $B$.

(3) Either $H\subset A$ or $H\subset B$, or both $H\cap hA$ and $H\cap hB$ are finite for all $h\in H$. Note that if $H\subset A$ or $H\subset B$, then $\mathcal T_1$ is a single vertex. 

(4) If $H\not\subset A$ and $H\not\subset B$, then either all cosets $hA$ in $\mathcal T_1$ are of valence 1 and all cosets $hB$ in $\mathcal T_1$ have the same finite valence which is equal to 1 plus the index of $H\cap C$ in $H\cap A$, or all cosets $hB$ in $\mathcal T_1$ are of valence 1 and all cosets $hA$ in $\mathcal T_1$ have the same finite valence which is equal to 1 plus the index of $H\cap C$ in $H\cap B$. Furthermore, if $hA$ is not valence 1, then all but one edge at $hA$ ends with an (valence 1) $h'B$ vertex for some $h'\in H$. Similarly if $hB$ is not valence 1. 

(5) Assume that $A$ is closer to $Q$ than is $B$ (so that $B$ has valence 1 in $\mathcal T_1$). If $h_1,h_2\in H$ and $h_1A\ne h_2A$ then there are no $hA$ or $hB$ vertices between $h_1A$ and $h_2A$ in $\mathcal T_1$. Similarly for $B$. 

(6) The vertex $Q$ is the midpoint of the geodesic between $A$ and $hA$ for some $h\in H$. 
\end{theorem}

\begin{proof}
Let $\Gamma$ be the Cayley graph of $G$ on generating sets for $A$, $B$ and $C$.
Call the vertices of $\mathcal T$ that contain elements of $H$, $H$-vertices. Note that if $h\in H$ is such that $h\in gA$ for some $g\in G$, then $gA=hA$. So the $H$-vertices of $\mathcal T$ are all cosets of the form $hA$ or $hB$ for $h\in H$. 
Since $H$ fixes $Q$, conclusion (1) is immediate. Conclusion (2) follows from (1) and the following fact: If $\alpha$ is a geodesic in $\mathcal T_1$ and $M$ is a closest point of $\alpha$ to $Q$ then the geodesic from $Q$ to $M$, followed by the segment of $\alpha$ from $M$ to an end point of $\alpha$, is a geodesic. This fact and (1) certainly imply that $M$ is the midpoint of $\alpha$. Note that $h_2h_1^{-1}$ maps $h_1A$ to $h_2A$ and so maps the geodesic from $Q$ to $h_1A$ to the geodesic from $Q$ to $h_2A$. Hence $h_2h_1^{-1}$ stabilizes the segment between $M$ and $Q$. Conclusion $(2)$ is verified. 

If $hA\cap H$ is infinite, for some $h\in H$, then for any $h'\in H$, $(h'(hA\cap H))=(hh'A)\cap H$ is infinite and so $hA\cap H$ is infinite for all $h\in H$. Suppose $A\cap H$ is infinite. If $hA\ne A$ for some $h\in H$ and $E$ is an edge in $\mathcal T$ separating $A$ and $hA$, then the finite subset of $\Gamma$ corresponding to $E$ separates infinitely many elements of $H$ in $A$ from infinitely many elements of $H$ in $hA$. But this implies that $H$ is not 1-ended in $G$. Instead, if $A\cap H$ is infinite, then $hA=A$ for all $h\in H$, and so $H\subset A$ and $\mathcal T_1$ consists of a single vertex. Conclusion (3) is verified. 

\medskip

$(\dagger)$ From this point on we assume that $hA\cap H$ and $hB\cap H$ are finite for all $h\in H$. 

\medskip 

Consider the vertex $A$ of $\mathcal T_1$ and an element $h\in H$ such that $hA\ne A$. %and there is no vertex $h'A$ on the geodesic between $A$ and $hA$ for $h'\in A$. %Let $E$ be the first edge of the geodesic between $A$ and $hA$. 
%Consider the vertices $HV_0$. If this is a single vertex, then $H$ stabilizes $V_0$ and we are finished. Assume $HV_0$ contains at least two vertices.  By the minimality of $\mathcal T_1$, every vertex of $\mathcal V_1$ is between two vertices of $HV_0$. Let $V_1\ne V_0$ be a closest vertex to $V_0$ such that $V_1\in HV_0$. 
Suppose $E$ is an edge of $\mathcal T_1$ separating $A$ and $hA$. If there are infinitely many distinct elements of $H$ contained in the $\mathcal T_1$ vertices (cosets) on the $A$ side of $E$ and infinitely many elements of $H$ contained in the $\mathcal T_1$ vertices on the $hA$ side of $E$, then the finite graph corresponding to $E$ in $\Gamma$ would separate two infinite sets of elements of $H$. But then $H$ would not be 1-ended in $\Gamma$. 

Instead, there are only finitely many elements of $H$ in the union of all $hA$ and $hB$ cosets on exactly one side of $E$. We are interested in the case where $A$ is a vertex of $E$. (See Figure \ref{FigEndPair1}). The last part of conclusion $(2)$ implies that $hE$ is the last edge of the geodesic between $A$ and $hA$. 

\begin{figure}
\vbox to 3in{\vspace {-2in} \hspace {-.1in}
\hspace{-1 in}

\vss }
\vspace{-3in}
\caption{Separating elements of $H$} 
\label{FigEndPair1}
\end{figure}

Suppose there are infinitely many elements of $H$ on the $A$ side of $E$ and say there are $K$ elements of $H$ on the $hA$ side of $E$.  Now consider the graph isomorphism $h$ acting on $\mathcal T_1$. %If $hE$ is not the last edge of the geodesic between $A$ and $hA$ then there are $K$ elements of $H$ on the $h^2A$ side of $hE$.  All of these elements of $H$ as well as $h$ are on the $hA$ side of $E$. But that is impossible since there are only $K$ such elements on the $hA$ side of $E$. 
Since $hE$ is the last edge of the geodesic between $A$ and $hA$, then there are infinitely many elements of $H$ on $hA$ side of $hE$, but all of these elements of $H$ are on the $hA$ side of $E$. This contradicts the fact that there are only $K$ elements of $H$ on the $hA$ side of $E$. 
This implies: 

\medskip

$(I)$ Suppose $h\in H$, $A\ne hA$ and $E$ is the first edge of a geodesic in $\mathcal T_1$ from $A$ and $hA$, then there are only finitely many elements of $H$ on the $A$ side of $E$ in $\mathcal T_1$, and infinitely many elements of $H$ on the $hA$ side of $E$.  
Similarly for $B$. 
 
%$(II)$ Under the hypotheses of $(I)$, $hE$ is the last edge of the geodesic from $A$ to $hA$. Similarly for $B$. 

\medskip

%Say there are $K$ elements of $H$ on the $A$ side of $E$. Suppose $hE$ is not the last edge of the geodesic from $A$ to $hA$. By translation, there are $K$ elements of $H$ on the $hA$ side of $hE$. But the $K$ elements of $H$ on the $A$ side of $E$ as well as $h$ are on the $hA$ side of of $hE$ (the desired  contradiction). 
Next we show:

\medskip

$(II)$ Under the hypotheses of $(I)$,  if $E'\ne E$ is an edge containing $A$, then the branch of $\mathcal T_1$ on the side of $E'$ opposite $A$ cannot contain a vertex $h'A$ for any $h'\in H$. Similarly for $B$. 

\medskip

Otherwise, $(I)$ implies that there are only finitely many elements of $H$ on the $A$ side of $E'$. (See Figure \ref{FigEndPair1}.) But all  (infinitely many) elements of $H$ on the $hA$ side of $E$ are on the $A$ side of $E'$, a contradiction and $(II)$ is verified.  


%Let $D$ be the first edge of the geodesic between $B$ and $hB$. We have:

%\medskip

%$(\ast\ast )'$ If $D'\ne D$ is an edge containing $B$, then the branch of $\mathcal T_1$ on the side of $D'$ opposite $B$ cannot contain a vertex $h'B$ for any $h'\in H$. 

\medskip

Again we assume the hypotheses of $(I)$: Suppose $h\in H$, $A\ne hA$ and $E$ is the first edge of a geodesic $\alpha$ between $A$ and $hA$. Let $F$ be the edge connecting $A$ and $B$ in $\mathcal T$. Then $F$ is an edge of $\mathcal T_1$. 
If $F\ne E$, then $(II)$ (with $F$ playing the role of $E'$) implies:

\medskip

\noindent $(\ast)$ There is no $h'\in H$ such that $h'A$ is a vertex of $\mathcal T_1$ on the side of $F$ opposite $A$. 

\medskip

\begin{figure}
\vbox to 3in{\vspace {-1.5in} \hspace {.6in}
\hspace{-1 in}

\vss }
\vspace{-1.7in}
\caption{Separating elements of $H$} 
\label{FigEndPair2}
\end{figure}

\noindent Since $hB$ is adjacent to $hA$ we have that $B\ne hB$. We have the configuration of Figure \ref{FigEndPair2} (see the last part of conclusion $(2)$). We apply $(II)$ to $B$ and $F$ to see that the branch of $\mathcal T_1$ on the side of $F$ opposite $B$ cannot contain a vertex $h'B$ for any $h'\in H$. Combining with $(\ast)$, we have that if $F\ne E$, then $B$ has valence 1 in $\mathcal T_1$.  Note that if $F=E$ then $A$ is valence 1 in $\mathcal T_1$. In any case: 

\medskip

$(III)$ Under the hypotheses of $(I)$, either $A$ is on the geodesic between $B$ and $hB$ and $B$ has valence 1 in $\mathcal T_1$ or $B$ is on the geodesic between $A$ and $hA$ and $A$ has valence 1 in $\mathcal T_1$. 

Without loss, assume that $B$ (and hence $h'B$ for all $h'\in H$) has valence 1 in $\mathcal T_1$.

\medskip

Next suppose $E'$ is an edge containing $A$, $F\ne E'\ne E$ and $h_1B$ is on the side of $E'$ opposite $A$. Then $h_1A$ is adjacent to $h_1B$ and $(II)$ implies that $h_1A=A$. Then $h_1\in A$. Since $H\cap A$ is finite, we have: 

\medskip

$(IV)$ The valence of $A$ in $\mathcal T_1$ is finite and with the exception of $E$, each edge $E'$ of $\mathcal T_1$ containing $A$ has as other vertex $h'B$ for some $h'\in H\cap A$. Observe that for the number of $H\cap C$ cosets in  $H\cap A$ is the number of edges at the vertex $A$ in $\mathcal T_1$ with end point $hB$ for some $h\in H$. 

\medskip

This completes the proof of conclusion (4).

Since $h'B$ has valence 1 for all $h'\in H$, there is no such vertex in the geodesic between $A$ and $hA$. Statement $(II)$ implies there is no vertex $h'A$ (with $h'\in H$), between $A$ and $hA$. Together this implies conclusion conclusion (5). 

For each $h\in H$ let $M_h$ be the midpoint of the geodesic between $A$ and $hA$. Choose $h\in H$ such that the distance between $M_h$ and $Q$ is minimal over all $h\in H$. If $h'\in H$, then  
Conclusion $(2)$ implies that $M_h$ and $M_h'$ are on the geodesic from  $Q$ to $A$. By the minimality of $h$, $M_{h}$ is between $M_{h'}$ and $Q$. Conclusion $(2)$ implies that $h'$ stabilizes $M_{h'}$ and $Q$ and hence $h'$ stabilizes $M_h$ for all $h'\in H$. This implies that $M_h=Q$. 
\end{proof}

\end{document}
